\documentclass[a4paper,11pt]{article}
\usepackage[utf8]{inputenc}
\usepackage[french]{babel}
\usepackage[T1]{fontenc}
\usepackage{amsmath,amssymb,amsthm}
\usepackage{geometry}
\geometry{margin=2cm}
\usepackage{enumitem}
\usepackage{hyperref}
\usepackage{booktabs}
\usepackage{xcolor}
\usepackage{listings}
\usepackage{titlesec}
\usepackage{fancyhdr}
\usepackage{lastpage}
\AtBeginDocument{\shorthandoff{;:!?}}

\titleformat{\section}{\large\bfseries\color{blue!60!black}}{}{0pt}{}[\vspace{-0.5ex}\rule{\textwidth}{0.4pt}]
\titleformat{\subsection}{\bfseries\color{blue!40!black}}{}{0pt}{}

\lstset{basicstyle=\ttfamily\small,breaklines=true,frame=single,aboveskip=6pt,belowskip=6pt}

\newtheorem*{important}{\`A retenir}
\pagestyle{fancy}
\fancyhf{}
\fancyhead[L]{\small STA211 -- Compl\'ement Tanagra TF/Keras R}
\fancyhead[R]{\small Fiche r\'ecapitulative}
\fancyfoot[C]{\thepage/\pageref{LastPage}}
\renewcommand{\headrulewidth}{0.4pt}

\title{\textbf{Fiche r\'ecapitulative -- STA211}\\[4pt]
\large Compl\'ement \texttt{fr\_Tanagra\_Tensorflow\_Keras\_R} \\[2pt]
\small TensorFlow / Keras sous R -- Perceptron simple et multicouche}
\author{Tanagra Data Mining -- R. Rakotomalala (avril 2018)}
\date{14 juin 2026}

\begin{document}
\maketitle
\thispagestyle{empty}

\begin{abstract}
Ce tutoriel pr\'esente l'utilisation des librairies TensorFlow / Keras sous R,
en suivant les m\^emes \'etapes que le document \'equivalent sous Python.
La transposition R $\leftrightarrow$ Python est particuli\`erement simple.
\end{abstract}

\vfill
\tableofcontents
\newpage

% ===================================================================
\section{Installation sous R}
% ===================================================================

\begin{lstlisting}
install.packages("keras")
library(keras)
install_keras()  # installe Keras + TensorFlow (backend)
\end{lstlisting}

\noindent\texttt{install\_keras()} installe \`a la fois le c\oe ur de Keras et le moteur
TensorFlow. Keras n'est qu'un \emph{front-end} : c'est TensorFlow qui effectue
les calculs.

% ===================================================================
\section{Donn\'ees}
% ===================================================================

Probl\`eme de discrimination binaire dans le plan. 2000 observations,
2 variables pr\'edicatives $X_1, X_2$. Fronti\`ere parabolique~:

\[
\text{Si } 0.1 \times X_2 > X_1^2 \implies \text{pos}, \text{ sinon neg}
\]

Pr\'eparation~:
\begin{lstlisting}
D <- read.table("artificial2d_data2.txt", sep="\t", header=T, dec=".")
y <- ifelse(D$Y=="pos", 1, 0)
set.seed(100)
idTrain <- sample(1:nrow(D), 1500, replace=F)
XTrain <- as.matrix(D[idTrain, 1:2])
yTrain <- y[idTrain]
XTest  <- as.matrix(D[-idTrain, 1:2])
yTest  <- y[-idTrain]
\end{lstlisting}

\noindent 1500 observations en apprentissage, 500 en test.

% ===================================================================
\section{Perceptron simple}
% ===================================================================

Architecture~:
\begin{lstlisting}
modelSimple <- keras_model_sequential()
modelSimple %>%
  layer_dense(units=1, input_shape=c(2), activation="sigmoid")
\end{lstlisting}

\begin{itemize}
  \item 1 neurone de sortie (codage 1/0), sigmo\"ide.
  \item 3 param\`etres \`a estimer : $a_0$ (biais) $+ a_1 X_1 + a_2 X_2$.
  \item Biais activ\'e par d\'efaut (\texttt{use\_bias = True}).
\end{itemize}

Compilation et apprentissage~:
\begin{lstlisting}
modelSimple %>% compile(
  loss="binary_crossentropy",
  optimizer="adam",
  metrics="accuracy"
)
modelSimple %>% fit(x=XTrain, y=yTrain, epochs=150, batch_size=10)
\end{lstlisting}

Poids estim\'es~: $a_0 = -1.283$, $a_1 = 0.256$, $a_2 = 1.947$.

\paragraph{Performances~:}
\begin{itemize}
  \item Apprentissage~: loss = 0.6429, accuracy~=~62.2\,\%
  \item Test (predict\_class + confusionMatrix)~: accuracy~=~67.8\,\%
  \item \'Evaluation Keras (\texttt{evaluate})~: loss~=~0.6206, acc~=~67.8\,\%
\end{itemize}

% ===================================================================
\section{Perceptron multicouche}
% ===================================================================

Architecture~: 1 couche cach\'ee (3 neurones) + 1 couche de sortie (1 neurone),
activation sigmo\"ide partout.

\begin{lstlisting}
modelMc <- keras_model_sequential()
modelMc %>%
  layer_dense(units=3, input_shape=c(2), activation="sigmoid") %>%
  layer_dense(units=1, activation="sigmoid")
\end{lstlisting}

Nombre de param\`etres~:
\begin{itemize}
  \item Entr\'ee $\rightarrow$ cach\'ee~: $(2 \times 3) + (1 \times 3) = 9$
  \item Cach\'ee $\rightarrow$ sortie~: $(3 \times 1) + (1 \times 1) = 4$
  \item \textbf{Total~: 13 param\`etres}
\end{itemize}

Compilation et apprentissage~: m\^emes param\`etres (binary\_crossentropy,
adam, 150 epochs, batch\_size=10).

\paragraph{Performances~:}
\begin{itemize}
  \item \textbf{Accuracy test~: 97.4\,\%}
  \item Nette am\'elioration par rapport au perceptron simple (67.8\,\%).
  \item Le r\'eseau non lin\'eaire (couche cach\'ee) capture mieux la
    fronti\`ere parabolique.
  \item Taux de reconnaissance sur l'\'echantillon test.
\end{itemize}

% ===================================================================
\section{Comparaison R vs Python}
% ===================================================================

\begin{center}
\begin{tabular}{lcc}
\toprule
Aspect & R & Python \\
\midrule
Syntaxe & \texttt{\%\textgreater{}\%} & API fonctionnelle \\
Installation Keras & \texttt{install\_keras()} & \texttt{pip install keras} \\
Fonctions & \texttt{layer\_dense()} & \texttt{layers.Dense()} \\
\bottomrule
\end{tabular}
\end{center}

\begin{important}
  Les modes op\'eratoires sont tr\`es similaires sous R et Python.
  Passer d'un langage \`a l'autre n'est pas un probl\`eme.
\end{important}

% ===================================================================
\newpage
\section*{Questions cl\'es (QCM)}
\addcontentsline{toc}{section}{Questions cl\'es (QCM)}

\begin{enumerate}
  \item Quelle fonction installe Keras + TensorFlow sous R~?
    \textbf{R. \texttt{install\_keras()}.}
  \item Quel op\'erateur cha\^ine les instructions Keras sous R~?
    \textbf{R. \texttt{\%\textgreater{}\%} (pipe).}
  \item Combien de param\`etres pour le perceptron simple~?
    \textbf{R. 3 (biais + 2 poids).}
  \item Combien de param\`etres pour le MLP (3 neurones cach\'es)~?
    \textbf{R. 13 (9 + 4).}
  \item Quelle est l'accuracy du perceptron simple en test~?
    \textbf{R. 67.8\,\%.}
  \item Quelle est l'accuracy du MLP en test~?
    \textbf{R. 97.4\,\%.}
  \item Quelle fonction d'activation est utilis\'ee~?
    \textbf{R. Sigmo\"ide (sigmoid).}
  \item Quel optimiseur est utilis\'e~?
    \textbf{R. Adam.}
  \item Quelle fonction de perte (loss)~?
    \textbf{R. \texttt{binary\_crossentropy}.}
  \item Quel est le nom de l'instruction de cr\'eation du mod\`ele~?
    \textbf{R. \texttt{keras\_model\_sequential()}.}
\end{enumerate}

\vfill

\end{document}
